\textbf{COLUNI 2023 -- Questão 37}\par\medskip Analise os documentos abaixo.

\par

\textbf{Documento I} “14 de setembro [de 1941]. Os alemães atiraram em nós com canhões de longo alcance. Ontem nosso bairro ficou à mercê dos bombardeios. Mesmo assim, nosso edifício permaneceu intacto – mas ao redor, literalmente ao redor, caíram obuses: na rua Ivanovskaia, na rua Raziezjaia, no n. 16, na pracinha Vladimir, na rua Marat, na rua Pravda, perto do Grande Teatro de Arte Dramática, a pouca distância do Teatro Alexandrinski, e em outros pontos mais. Os obuses passaram por cima de nosso edifício sem nos atingir, mas caíram perto. A qualquer momento podemos ser mortos. [...] Todo mundo quer viver, é óbvio! E os que foram mortos também tinham vontade de viver. Dentre esses mortos há crianças, bebês, velhos e jovens: moças e rapazes querendo viver. [...]

\par

Nós de Leningrado, durante esses terríveis dias 8, 9 e 10 de setembro, nos habituamos ao fato de que sempre, por volta das 11h, a sirene começa a tocar e corremos para os abrigos (aqueles, em todo caso, para quem a vida tem um preço). O espetáculo então começa: mísseis, trovoar dos aviões, estrondos das bombas explosivas, assobio das bombas incendiárias [...]”

\par

Obus: projétil onde se colocam explosivos que são lançados em caso de guerra.

\par

Lena Mukhina. O diário de Lena: a história real de uma adolescente durante a Segunda Guerra. In: CHIBA, Charles Hokiti Fukushigue, MINORELLI, Caroline Torres. \textbf{Convergências}: história, 9º ano. São Paulo: Edições SM, 2016. p. 88.

\par

\textbf{Documento I}I 13 de janeiro de 1943. “[...] Coisas terríveis estão acontecendo lá fora. A qualquer hora do dia ou da noite pessoas pobres e desamparadas são retiradas de suas casas. Não têm permissão de levar nem mesmo uma sacola com alguma coisa e um pouco de dinheiro, e, mesmo quando têm, essas posses lhes são roubadas no caminho. Famílias são rompidas; homens, mulheres e crianças são separados. Crianças chegam da escola e descobrem que os pais desapareceram. Mulheres voltam das compras e descobrem as casas lacradas e que as famílias desapareceram. Os cristãos holandeses também estão com medo porque seus filhos são mandados à Alemanha. Todo mundo anda apavorado. Todas as noites centenas de aviões passam sobre a Holanda a caminho das cidades alemãs, para semear suas bombas em solo alemão. Toda hora centenas, ou talvez milhares, de pessoas são mortas na Rússia e na África. Ninguém pode ficar longe do conflito, o mundo inteiro está em guerra, e, mesmo com os Aliados se saindo melhor, o fim não está próximo. [...]

\par

As crianças deste bairro andam com camisas finas e sapatos de madeira. Não têm casacos, nem capas, nem meias, nem ninguém para ajudá-las. Mordendo uma cenoura para acalmar as dores da fome, saem de suas casas frias e andam pelas ruas até salas de aula ainda mais frias. As coisas ficaram tão ruins na Holanda que hordas de crianças abordam os pedestres para implorar um pedaço de pão. [...]”

\par

Horda: grupo numeroso e desorganizado de pessoas.

\par

Anne Frank. O diário de Anne Frank. In: CHIBA, Charles Hokiti Fukushigue, MINORELLI, Caroline Torres. \textbf{Convergências}: história, 9º ano. São Paulo: Edições SM, 2016. p. 89.

\par

Sobre os diários e os temas por eles tratados nos documentos acima, é INCORRETO afirmar:\begin{enumerate}
\item Os diários pessoais são fontes históricas preciosas para compreendermos a vida cotidiana das pessoas
        que viveram em determinados contextos.
\item Os diários, escritos pelas adolescentes Lena Mukhina e Anne Frank, retratam a violência militar, física e
        psicológica sofrida por crianças, jovens, velhos e adultos que moravam em diferentes cidades, durante a
        Segunda Guerra Mundial.
\item Os diários podem ser compreendidos como documentos que expressavam os sentimentos das
        adolescentes em relação à guerra e à dor que ela causava.
\item Os diários nos permitem identificar os instrumentos utilizados na guerra e compreender os efeitos, para os
        povos russo e judeu, do desembarque aliado na Normandia.\end{enumerate}
